%% kapa: exam=1 round=30 year=2019 period=2 form=A subject=appraisal_law subject_ko=감정평가관계법규 seq=10 number=10
%% source: Q-Net 감정평가사 기출문제 (cst00309) artlSeq=5205205 file=extracted/1008/5205205/2019년 제30회 감정평가사 1차시험 문제지/2019년 제30회 감정평가사 1차시험 문제지/2019년도 제30회 감정평가사 1차 2교시 A형.hwp
%% conversion: kordoc 4.24.51 -> tools/gen1.py (auto); review=proofread-sonnet; flags=none
\begin{question}{10}
국토의 계획 및 이용에 관한 법령상 공동구에 관한 설명으로 옳지 않은 것은?
\begin{choices}
  \item 「택지개발촉진법」에 따른 100만 제곱미터의 택지개발지구에서 개발사업을 시행하는 자는 공동구를 설치하여야 한다.
  \item 공동구관리자는 5년마다 해당 공동구의 안전 및 유지관리계획을 수립ㆍ시행하여야 한다.
  \item 공동구의 관리에 소요되는 비용은 그 공동구를 점용하는 자가 함께 부담하되, 부담비율은 점용면적을 고려하여 공동구관리자가 정한다.
  \item 공동구 설치비용 부담액을 완납하지 않은 자가 공동구를 점용하려면 그 공동구를 관리하는 공동구관리자의 허가를 받아야 한다.
  \item 공동구관리자는 1년에 1회 이상 공동구의 안전점검을 실시하여야 한다.
\end{choices}
\end{question}
\answer{1}
